% 第 18 章 Working With Manufacturing Objects
\bichapter{处理制造对象}{Working With Manufacturing Objects}
\label{chap:gui-mfg}

IC Compiler II GUI 可以查看、选择和编辑设计中的下列制造对象。它们参与半导体芯片的制造，因此称为制造对象。
\begin{itemize}
  \item Assembly die：从晶圆上切下之后的物理裸片。
  \item Sealring：比芯片边界更大的矩形晶圆区域，用于保护裸片免受机械损伤。
  \item Stitching zone（Kerf）：把晶圆切开形成裸片时被去除的材料宽度。
\end{itemize}

相关主题：
\begin{itemize}
  \item 选择制造对象
  \item 查询制造对象
  \item 查看制造对象的属性
  \item 编辑制造对象
\end{itemize}

\section{选择制造对象}
\subsection*{Selecting a Manufacturing Object}

可以用下列方式选择制造对象：
\begin{itemize}
  \item 在布局视图中单击该对象；
  \item 在上下文菜单中选择 Related Manufacturing Shapes；
  \item 使用 Select by Name 工具栏。
\end{itemize}

指针停在制造对象上时，工具提示会显示该对象的信息。在布局视图中，用 View Settings 面板控制制造对象的可见性与可选性。

\section{查询制造对象}
\subsection*{Querying Manufacturing Objects}

用 Query 面板搜索制造对象。

\section{查看制造对象的属性}
\subsection*{Viewing the Properties of a Manufacturing Object}

用 Properties 面板查看制造对象的属性。

\section{编辑制造对象}
\subsection*{Editing a Manufacturing Object}

可以在 GUI 中拉伸、添加或切割制造对象。
\begin{enumerate}
  \item 在布局视图中选中对象。
  \item 使用 Edit 菜单中的编辑工具。
\end{enumerate}

在布局视图中编辑对象的方法，见第~\ref{chap:gui-start}~章「编辑物理版图」。
